\subsection{A senior dividend and a cap change who receives the next unit of profit}
\label{capped-participation-eps}

The proportional example leaves every additional unit of earnings to the same
sharing rule. A senior entitlement and a cap change that allocation. Consider
100 ordinary shares and 20 nonconvertible participating preference instruments,
all outstanding throughout the year. The supplied terms give the preference
class a cumulative current-year entitlement of 20, then equal participation
per instrument in positive residual earnings, up to total class earnings of
40. The ordinary shares receive what remains. The preference class does not
share residual losses. Equity classification and these rights are assumed
facts of this example; the earnings allocation does not declare a dividend or
create a payable. IAS~33.14(b) includes the current period's cumulative
entitlement whether declared or not, while A13--A14 require allocation by the
actual participation rights \citep[IAS~33.14(b), A13--A14]{eu20231803}.

At allocable profit 140, deduct the senior 20 first. The remaining 120 divides
over 120 equal participation units, so the preference class receives another
20 and reaches its cap. Ordinary earnings are 100 and EPS is 1. At profit 80,
the residual is only 60: the preference class receives 20 plus 10, leaving
50 for ordinary shareholders and EPS of one half. At profit 200, an uncapped
residual allocation would give the preference class another 30. Its cap allows
only 20, leaving ordinary earnings of 160. The ordinary shareholders receive
every further unit once the cap binds.

To state that rule precisely, let \(E\) be earnings available to these two
classes after all other required attribution, \(n\) the ordinary population,
and \(m\) the total residual participation weight of the preference class.
Let \(f\) be its senior period entitlement and \(c\) the maximum additional
participation. Thus this example has \(m=20\), \(f=20\) and \(c=20\);
the total cap is \(f+c=40\). The rights give
\[
 p=\min\!\left(c,\frac{m\max(E-f,0)}{n+m}\right),\qquad
 E_P=f+p,\qquad E_O=E-f-p,\qquad r=\frac{E_O}{n}.
\]
Here \(p\) is residual participation, \(E_P\) and \(E_O\) are the two
class allocations, and \(r\) is ordinary EPS. The two allocations sum to
\(E\). If earnings are zero, the current cumulative entitlement still assigns
20 to the preference class and \(-20\) to the ordinary class; ordinary EPS
is \(-1/5\). At earnings of \(-40\), ordinary earnings are \(-60\), giving
EPS \(-3/5\). These negative ordinary amounts follow from the senior right
and the absence of residual loss sharing, rather than a rule that every class
must receive a nonnegative allocation.

Assumed issuance now changes both sides of the ordinary EPS fraction. Return
to profit 140 and add 100 options with exercise price 10 and average market
price 20. The IAS~33.45--47 treasury-share calculation gives
\(100-100(10)/20=50\) incremental ordinary shares
\citep[IAS~33.45--47]{eu20231803}. Assume no further service proceeds, no
earnings effect and a full-year eligible option population. A14 includes the
assumed shares in the ordinary population. With \(n=150\), the preference
class's residual participation falls to \(120(20)/170=240/17\), below its
cap. Consequently
\[
 E_P=20+\frac{240}{17}=\frac{580}{17},\qquad
 E_O=140-\frac{580}{17}=\frac{1800}{17},\qquad
 r=\frac{1800/17}{150}=\frac{12}{17}.
\]
The preference class receives less because more ordinary shares now share the
residual. Keeping ordinary earnings at 100 would instead give \(2/3\), which
is wrong for these participation rights. The numerator cannot be frozen merely
because the option has no direct earnings adjustment.

Add a full-year convertible debt series whose complete after-tax earnings
effect is 6 and whose incremental ordinary population is 10. With both series
assumed, total earnings become 146 and the ordinary population becomes 160.
The preference class receives \(20+126(20)/180=34\); ordinary earnings are
112 and EPS is \(112/160=7/10\). This is below \(12/17\), so the debt
remains dilutive after the option. Debt alone would give ordinary EPS
\((146-20)/(110+20)=63/65\). The same conversion has different ordinary
numerator effects depending on the shares already assumed.

There is nevertheless a simple ordering rule for this restricted family.
Assume continuing and total earnings coincide, rights remain fixed, and each
series independently adds a complete nonnegative gross after-tax earnings
amount \(a_i\) and a positive ordinary share amount \(b_i\). The gross
amount is measured before the two-class allocation. For \(E\geq f\),
substituting the participation formula gives
\[
 r(E,n)=\max\!\left(\frac{E-f}{n+m},\frac{E-f-c}{n}\right).
\]
The first fraction applies before the cap binds; the second applies once the
preference class receives its full cap. An assumed series adds the same
\(a_i\) to either numerator and \(b_i\) to either denominator. If current
EPS is \(r\geq0\), both resulting fractions fall below \(r\) exactly when
\(a_i-rb_i<0\). If that expression is positive, the previously active
fraction rises; equality leaves EPS unchanged. Thus order complete gross
effects by \(a_i/b_i\), recalculate the allocation, and include a series only
when it strictly lowers ordinary EPS. IAS~33.41--44 provides the continuing-
operations control and separate-series dilution objective
\citep[IAS~33.33--44]{eu20231803}; the formula here is a local derivation for
the stated rights and additive gross effects.

The ordering also finds the minimum over this family's combinations. For any
trial bound \(r\), requiring both fractions to be at most that bound imposes
two inequalities with the same varying sum
\(\sum_i(a_i-rb_i)\). Including every negative term minimizes that sum in
both inequalities, so a minimizing combination exists among the prefixes in
ascending \(a_i/b_i\) order. An included series leaves EPS above its own
ratio, and all later ratios are at least as large. Once a series no longer
lowers EPS, later ones cannot restore dilution. Below the senior entitlement,
nonnegative earnings additions and positive share additions cannot deepen the
ordinary loss per share. Equality is excluded by the strict dilution condition.

The read-only \texttt{CappedParticipationEPS/1} method uses exact rational
amounts and returns the allocation before and after each trial assumption.
It requires explicit rights, complete effects, context and period findings;
unknown findings leave the answer unresolved, and unsupported or contradictory
terms produce no partial EPS result. Its 18 authored cases include both cap
regimes, zero and negative earnings, equality, reversed input order and the
nonparticipating limit. An independent calculation enumerates all 49 finite
exercise combinations. This supports the specified calculation, while the
existing fixed ordinary-numerator method retains its separate assumptions.

Complete gross effects are the consequential restriction. IAS~33.35 points
out that removing interest expense can increase a nondiscretionary employee
profit-sharing expense. If that expense depends on the combined conversions,
single-series effects need not add. Conditional exercise or changed rights
can alter the possible assumptions as well. Those cases need their own joint
earnings and eligibility calculation before any dilution ordering is justified.
